% Section 8 -- Conclusion and Future Work
% Target: 1 page.

DLC unifies authorization logic, symbolic cryptography, computational
cryptography, runtime wire format, and reproducible CI in a single
artifact where the same chain-splicing-defeating rule appears at every
layer.  The unification is the contribution; the individual pieces
build on Garg-Pfenning (logic), Cremers-Blanchet (symbolic crypto),
and Tuma et al.~(VCVio, computational crypto in Lean).

\paragraph{What is proven.}  Twelve theorems with no \texttt{sorry}
and no new axioms:
\begin{itemize}
\item T4 non-introduction direction (\texttt{pendingObligations\_shift}
  closed; substitution-subset and step-case-analysis open).
\item Three graded comonad laws (identity, associativity, DpBudget
  saturating-add associativity).
\item Four strong indexed monad laws (left-id, right-id,
  associativity, strength left-unit).
\item Four structural substitution lemmas (\texttt{shift\_zero} plus three
  \texttt{substAt\_var\_*} cases).
\item Three Tamarin lemmas (secrecy, non-splicing, executability).
\item Three ProVerif queries (same, by independent prover).
\end{itemize}

\paragraph{What is stated but open.}  T1 full-calculus decidability,
T2 cryptographic correspondence, T3 non-interference, T4 multiset
form, substitution composition, subject reduction, the \S 4.4
protocol-logic correspondence.  Each has a documented closure path:
T1 via porting the Rust \texttt{decide\_pure} to Lean; T2 via VCVio
EUF-CMA reduction; T3 via hand-rolled logical relations in Lean
without Iris; T4 multiset form via the M3 runtime layer's
\texttt{Discharged} evidence; substitution composition via Ramos
et al.'s 709-line Metatheory framework once that becomes available.

\paragraph{What is not in this paper.}  We make no claim to having
established DLC as a paradigm.  That bar requires (a) full closure of
the open theorems, (b) external citation by adjacent work,
(c) standards-track adoption, (d) deployment beyond the reference
implementation.  Each of these depends on actors beyond the authors
and on multi-quarter time horizons.  We position the present paper as
a \emph{credible foundation} for that future work, with the artifact's
double-prover symbolic verification and reproducible ledger as
methodological contributions on their own.

\paragraph{Future work.}  In the medium term, four closure milestones
in priority order:

\begin{enumerate}
\item Close T4 non-introduction in full (substitution-subset lemma,
  step-case analysis).  Autonomously reachable in $\sim$2 weeks.
\item Close T1 propositional soundness via a Lean \texttt{decideLean}
  mirroring the Rust implementation.  Autonomously reachable in
  $\sim$3 weeks.
\item Close T2 via VCVio.  Autonomously reachable in $\sim$5-8 weeks.
\item Close T3 atomic-fragment + \texttt{lift} cases as a logical
  relations skeleton.  Autonomously reachable in $\sim$5 weeks;
  full T3 needs $\sim$3-5 months including the Lean infrastructure
  layer.
\end{enumerate}

In the longer term, collaborator engagement: Pfenning/Garg on T1/T3
proof technique, Cremers on Tamarin/ProVerif review, Blanchet on the
VCVio EUF-CMA reduction's tightness, Myers on IFC and obligation
algebra.  An IETF SECDISPATCH presentation of the wire-format draft
is the natural standards-track entry point.

\paragraph{The diagonal as the contribution.}  We close by stating
the contribution claim in its strongest form: DLC is the first
delegation calculus in which the same chain-splicing-defeating rule
appears across (a) a Lean-mechanized type system, (b) a Tamarin
symbolic model, (c) a ProVerif independent symbolic model, (d) an
EasyCrypt computational skeleton, (e) a Rust runtime with content-
addressed wire format.  Each layer is independently verifiable; their
agreement is the artifact's load-bearing methodological claim.  Whether
DLC ultimately earns the name of a paradigm depends on closures and
adoption that lie beyond this paper's horizon.  What this paper
contributes is the diagonal itself, with enough mechanical evidence
along each axis to make the unification credibly defensible under
external review.
